% !TeX program = pdflatex
\documentclass[11pt,a4paper]{amsart}

\newcommand{\notelanguage}{finnish}
% Shared preamble for course notes and exercise sets.

\usepackage[T1]{fontenc}
\usepackage{lmodern}
\providecommand{\notelanguage}{english}
\usepackage[finnish,english]{babel}
\AtBeginDocument{\expandafter\selectlanguage\expandafter{\notelanguage}}

\newcommand{\notesfinnishlanguage}{finnish}
\ifx\notelanguage\notesfinnishlanguage
\newcommand{\notestheoremname}{Lause}
\newcommand{\noteslemmaname}{Lemma}
\newcommand{\notespropositionname}{Propositio}
\newcommand{\notescorollaryname}{Korollaari}
\newcommand{\notesdefinitionname}{Määritelmä}
\newcommand{\notesexamplename}{Esimerkki}
\newcommand{\notesproblemname}{Tehtävä}
\newcommand{\notesexercisename}{Tehtävä}
\newcommand{\notesremarkname}{Huomautus}
\newcommand{\notessolutionname}{Ratkaisu}
\else
\newcommand{\notestheoremname}{Theorem}
\newcommand{\noteslemmaname}{Lemma}
\newcommand{\notespropositionname}{Proposition}
\newcommand{\notescorollaryname}{Corollary}
\newcommand{\notesdefinitionname}{Definition}
\newcommand{\notesexamplename}{Example}
\newcommand{\notesproblemname}{Problem}
\newcommand{\notesexercisename}{Exercise}
\newcommand{\notesremarkname}{Remark}
\newcommand{\notessolutionname}{Solution}
\fi

\usepackage[a4paper,margin=30mm]{geometry}
\usepackage{microtype}
\usepackage{graphicx}
\usepackage{enumitem}

\newcommand{\exerciseheading}[2]{%
  \par\noindent
  {\large\bfseries #1\par}
  \vspace{0.1em}
  \noindent{\large #2\par}
  \vspace{1.5em}
}

\usepackage{amsmath}
\usepackage{amssymb}
\usepackage{amsthm}
\usepackage{mathtools}
\usepackage{aliascnt}

% Row operations beside the affected row, e.g. \rowop{+2R_1}.
\newcommand{\rowop}[1]{\mathrel{\big|}\,#1}

% Linear systems: the argument gives the number of alignment pairs.
\newenvironment{linearsystem}[1]
{\left\{
  \begin{alignedat}{#1}}
    {
  \end{alignedat}\right.}

  \usepackage[hidelinks,pdfusetitle]{hyperref}

  \numberwithin{equation}{section}
  \setlist{itemsep=0.25em,topsep=0.5em}

  \theoremstyle{plain}
  \newtheorem{theorem}{\notestheoremname}[section]
  \newaliascnt{lemma}{theorem}
  \newtheorem{lemma}[lemma]{\noteslemmaname}
  \aliascntresetthe{lemma}
  \newaliascnt{proposition}{theorem}
  \newtheorem{proposition}[proposition]{\notespropositionname}
  \aliascntresetthe{proposition}
  \newaliascnt{corollary}{theorem}
  \newtheorem{corollary}[corollary]{\notescorollaryname}
  \aliascntresetthe{corollary}

  \theoremstyle{definition}
  \newaliascnt{definition}{theorem}
  \newtheorem{definition}[definition]{\notesdefinitionname}
  \aliascntresetthe{definition}
  \newaliascnt{example}{theorem}
  \newtheorem{example}[example]{\notesexamplename}
  \aliascntresetthe{example}
  \newaliascnt{problem}{theorem}
  \newtheorem{problem}[problem]{\notesproblemname}
  \aliascntresetthe{problem}
  \newtheorem{exercise}{\notesexercisename}

  \theoremstyle{remark}
  \newaliascnt{remark}{theorem}
  \newtheorem{remark}[remark]{\notesremarkname}
  \aliascntresetthe{remark}

  \newenvironment{solution}
  {
  \begin{proof}[\notessolutionname]}
    {
  \end{proof}}

  \usepackage[nameinlink,noabbrev]{cleveref}

  \crefname{theorem}{theorem}{theorems}
  \crefname{lemma}{lemma}{lemmas}
  \crefname{proposition}{proposition}{propositions}
  \crefname{corollary}{corollary}{corollaries}
  \crefname{definition}{definition}{definitions}
  \crefname{example}{example}{examples}
  \crefname{problem}{problem}{problems}
  \crefname{exercise}{exercise}{exercises}
  \crefname{remark}{remark}{remarks}

  \ifx\notelanguage\notesfinnishlanguage
  \crefname{theorem}{lause}{lauseet}
  \crefname{lemma}{lemma}{lemmat}
  \crefname{proposition}{propositio}{propositiot}
  \crefname{corollary}{korollaari}{korollaarit}
  \crefname{definition}{määritelmä}{määritelmät}
  \crefname{example}{esimerkki}{esimerkit}
  \crefname{problem}{tehtävä}{tehtävät}
  \crefname{exercise}{tehtävä}{tehtävät}
  \crefname{remark}{huomautus}{huomautukset}
  \fi

  \newcommand{\NN}{\mathbb{N}}
  \newcommand{\NNone}{\mathbb{N}_1}
  \newcommand{\ZZ}{\mathbb{Z}}
  \newcommand{\QQ}{\mathbb{Q}}
  \newcommand{\RR}{\mathbb{R}}
  \newcommand{\CC}{\mathbb{C}}
  \newcommand{\PP}{\mathbb{P}}

  \DeclareMathOperator{\im}{im}
  \DeclareMathOperator{\rank}{rank}
  \DeclareMathOperator{\syt}{syt}
  \DeclareMathOperator{\pyj}{pyj}
  \DeclareMathOperator{\dom}{dom}
  \DeclarePairedDelimiter{\abs}{\lvert}{\rvert}
  \DeclarePairedDelimiter{\norm}{\lVert}{\rVert}
  \newcommand{\set}[1]{\left\lbrace#1\right\rbrace}
  \newcommand{\maxset}[1]{\max\set{#1}}
  \newcommand{\minset}[1]{\min\set{#1}}
  \newcommand{\powerset}[1]{\mathcal{P}\!\left(#1\right)}


\title{Lineaarialgebra ja matriisilaskenta I\\Harjoitus 3}
\author{}
\date{}

\begin{document}

\exerciseheading{Lineaarialgebra ja matriisilaskenta I}{Harjoitus 3}

\begin{exercise}
  Olkoon \(A =
    \begin{bmatrix}
      1 & 0 \\
      -5 & 2
  \end{bmatrix}\).
  \begin{enumerate}[label=(\alph*), leftmargin=*]
    \item Määritä sellaiset alkeismatriisit \(E_1\) ja \(E_2\), että
      \(E_2E_1A = I\).
    \item Kirjoita \(A^{-1}\) kahden alkeismatriisin tulona.
    \item Kirjoita \(A\) kahden alkeismatriisin tulona.
  \end{enumerate}
\end{exercise}

\begin{exercise}
  Olkoon
  \[
    A =
    \begin{bmatrix}
      2 & 3 & -1 \\
      1 & -1 & 2 \\
      1 & 2 & -1
    \end{bmatrix}.
  \]
  Jos \(
    \begin{bmatrix} 2 \\ -2 \\ 3
  \end{bmatrix}\) on yhtälöryhmän
  \(Ax = b\) eräs ratkaisu, niin mikä on sen yleinen ratkaisu? Sovella
  ratkaisussasi Lemmaa 3.1.1. Mitä voit päätellä matriisin \(A\)
  säännöllisyydestä neliömatriisien kääntyvyyslauseen 2.5.1 tai Lauseen
  3.10.1 perusteella?
\end{exercise}

\begin{exercise}
  Osoita, että avaruuden \(\RR^{2\times 1}\) joukko
  \[
    V_0 = \left\{
      \begin{bmatrix} x_1 \\ x_2
      \end{bmatrix}
      \in \RR^{2\times 1}
      \colon x_1 \in \RR \text{ ja } x_2 = 0
    \right\}
  \]
  on avaruuden \(\RR^{2\times 1}\) aliavaruus. Miksi avaruuden
  \(\RR^{2\times 1}\) joukot
  \[
    \begin{aligned}
      V_1 &= \left\{
        \begin{bmatrix} x_1 \\ x_2
        \end{bmatrix}
        \in \RR^{2\times 1}
        \colon x_1 \in \RR \text{ ja } x_2 = 1
      \right\}
      \\[1ex]
      \text{ja}\qquad
      V_2 &= \left\{
        \begin{bmatrix} x_1 \\ x_2
        \end{bmatrix}
        \in \RR^{2\times 1}
        \colon x_1 \in \RR \text{ ja } x_2 = x_1^2
      \right\}
    \end{aligned}
  \]
  eivät ole avaruuden \(\RR^{2\times 1}\) aliavaruuksia?
\end{exercise}

\begin{exercise}
  Määritä avaruuden \(\RR^{2\times 4}\) matriisin
  \[
    A =
    \begin{bmatrix}
      1 & 1 & 1 & 0 \\
      2 & 1 & 0 & 1
    \end{bmatrix}
  \]
  nolla-avaruus \(\operatorname{Null}(A)\). Tutki lisäksi, onko
  \(\operatorname{Col}(A) = \RR^{2\times 1}\).
\end{exercise}

\begin{exercise}
  Olkoot
  \[
    v_1 =
    \begin{bmatrix} 0 \\ -2 \\ 2
    \end{bmatrix},
    \quad
    v_2 =
    \begin{bmatrix} 1 \\ 3 \\ -1
    \end{bmatrix},
    \quad
    w_1 =
    \begin{bmatrix} 3 \\ 1 \\ 5
    \end{bmatrix},
    \quad
    w_2 =
    \begin{bmatrix} 0 \\ 4 \\ 5
    \end{bmatrix},
    \quad
    w_3 =
    \begin{bmatrix} 0 \\ 0 \\ 0
    \end{bmatrix}
    \in \RR^{3\times 1}.
  \]
  Tutki, mitkä vektoreista \(w_1\), \(w_2\) ja \(w_3\) ovat vektoreiden
  \(v_1\) ja \(v_2\) lineaarikombinaatioita. Toisin sanoen, mitkä
  vektoreista \(w_1\), \(w_2\) ja \(w_3\) kuuluvat joukkoon
  \(\operatorname{Sp}(v_1,v_2)\)?
\end{exercise}

\begin{exercise}
  Tutki, ovatko
  \begin{enumerate}[label=(\alph*), leftmargin=*]
    \item \(v_1 =
        \begin{bmatrix} 2 \\ 1
      \end{bmatrix}\) ja
      \(v_2 =
        \begin{bmatrix} 3 \\ 2
      \end{bmatrix}\) lineaarisesti
      riippumattomia \(\RR^{2\times 1}\):ssa,
    \item \(w_1 =
        \begin{bmatrix} 2 \\ 1 \\ -2
      \end{bmatrix}\),
      \(w_2 =
        \begin{bmatrix} 3 \\ 2 \\ -2
      \end{bmatrix}\) ja
      \(w_3 =
        \begin{bmatrix} 2 \\ 2 \\ 0
      \end{bmatrix}\) lineaarisesti
      riippumattomia \(\RR^{3\times 1}\):ssa.
  \end{enumerate}
\end{exercise}

\end{document}
